\section{The turning-point iteration}
\label{sec:iteration}

From the first turning point the algorithm repeatedly (i) solves the reduced KKT system
for the affine segment \eqref{eq:segment}; (ii) finds the largest $\lambda$ below the
current one at which an \emph{event} occurs; (iii) flips the status of the single
coordinate responsible; and (iv) records the new turning point, stopping at
$\lambda = 0$.

\subsection{Solving the reduced KKT system by block elimination}
\label{sec:blockelim}

Write $C_\F$ for the active constraint rows restricted to the free set ($A_\F$ when
$\Sset = \varnothing$; \S\ref{sec:ineq}) and $c$ for their right-hand side. The
reduced system is the saddle-point problem
\begin{equation}
\label{eq:saddle}
\begin{bmatrix} \Sig_{\F\F} & C_\F\trans \\[2pt] C_\F & 0 \end{bmatrix}
\begin{bmatrix} w_\F \\ \nu \end{bmatrix}
=
\begin{bmatrix} r_1 \\ r_2 \end{bmatrix},
\end{equation}
solved for two right-hand sides at once, the ``$\alpha$ system''
$r_1^{\alpha} = -\Sig_{\F\B} w_\B$, $r_2^{\alpha} = c - C_\B w_\B$, and the
``$\beta$ system'' $r_1^{\beta} = \mu_\F$, $r_2^{\beta} = 0$. Rather than factor the
full saddle matrix, the code eliminates the free block. A single multi-column solve
against $\Sig_{\F\F}$ gives $Y = \Sig_{\F\F}\inv C_\F\trans$ and
$z^{\alpha,\beta} = \Sig_{\F\F}\inv r_1^{\alpha,\beta}$. The multipliers come from the
small Schur complement $S = C_\F Y$. With $n_\F$ free assets and $m_c = m + |\Sset|$ active rows,
$C_\F \in \R^{m_c\times n_\F}$, $Y \in \R^{n_\F\times m_c}$ and $S \in \R^{m_c\times m_c}$,
\begin{equation}
\label{eq:schur}
\nu^{\alpha} = S\inv\!\left(C_\F z^{\alpha} - r_2^{\alpha}\right),
\qquad
\nu^{\beta} = S\inv C_\F z^{\beta},
\end{equation}
and back-substitution gives
\begin{equation}
\label{eq:freesoln}
\alpha_\F = z^{\alpha} - Y\nu^{\alpha},
\qquad
\beta_\F = z^{\beta} - Y\nu^{\beta}.
\end{equation}
The dominant cost is the solve against $\Sig_{\F\F}$; for the plain budget $S$ is a
scalar. With $1 < m_c \le n_\F$ the rows are first replaced by an orthonormal basis of
their span, $C_\F\trans = QR$, and the system is solved with $Q$ in place of
$C_\F\trans$, the target $R^{-\mathsf{T}} r_2$, and the multipliers recovered by a
triangular solve with $R$. The condition number of $C_\F \Sig_{\F\F}\inv C_\F\trans$ can
reach $\operatorname{cond}(\Sig_{\F\F})\operatorname{cond}(C_\F)^2$, while
$Q\trans\Sig_{\F\F}\inv Q$ is no worse conditioned than $\Sig_{\F\F}$, so nearly
dependent active rows lose accuracy in proportion to $\operatorname{cond}(C_\F)$, the
sensitivity of the constraint data, and not to its square. Solving only for the free block keeps every step a dense solve of size
$|\F|$, where \citet{markowitz2021} assemble one large KKT system over all variables.

To decide when a blocked asset should re-enter, the code forms the reduced gradients
\begin{equation}
\label{eq:gamma}
\gamma = \Sig\,\alpha + C\trans \nu^{\alpha}, \qquad
\delta = \Sig\,\beta + C\trans \nu^{\beta} - \mu,
\end{equation}
where $\nu$ stacks the equality and active-row multipliers, so that $\gamma_i + \lambda\delta_i$ is (proportional to) the bound multiplier of a
blocked asset $i$ along the segment; the products $\Sig\alpha$, $\Sig\beta$ go through
the covariance operator of \S\ref{sec:operators}.

\subsection{Event detection}
\label{sec:events}

The segment stays optimal until one of four events changes the active set. Two act on
variables through the box bounds:
\begin{itemize}
\item \textbf{Free $\to$ blocked.} A free weight $w_i = \alpha_i + \lambda\beta_i$ meets
its upper bound at $\lambda_i = (\ub_i - \alpha_i)/\beta_i$ when $\beta_i < 0$, or its
lower bound at $\lambda_i = (\lb_i - \alpha_i)/\beta_i$ when $\beta_i > 0$.
\item \textbf{Blocked $\to$ free.} A blocked asset's multiplier changes sign at
$\lambda_i = -\gamma_i/\delta_i$, admissible at the upper bound with $\delta_i < 0$ or
the lower bound with $\delta_i > 0$.
\end{itemize}
The other two are the same moves on the rows of $G w \le h$, described in
\S\ref{sec:ineq}. The candidates form a matrix $L\in\R^{(n+p)\times 4}$, one row per
variable or inequality row, with $-\infty$ where an event cannot occur, under two
numerical guards:
\begin{itemize}
\item \textbf{Slope floor.} A slope at round-off level is treated as zero. The floor is
$\sqrt{\varepsilon_{\text{mach}}}$ in the problem's own units, so the trace is
invariant to rescaling $\mu$ or $\Sig$ (Appendix~\ref{sec:guards}).
\item \textbf{$\lambda$ window.} Candidates above the current $\lambda$, beyond a
relative $10^{-10}$, lie outside the segment and are discarded; the same window decides
which events are tied.
\end{itemize}
The next turning point is $\lambda^{\star} = \max_{i,t} L_{it}$; if
$\lambda^{\star} \le 0$ the frontier is complete.

\subsection{General linear inequalities}
\label{sec:ineq}

A row of $G w \le h$ is the analogue of a box bound: either \emph{active}
($g_j\trans w = h_j$) or slack. While active it is carried through the reduced solve as
an extra equality row, stacked under $A$ into $C_\F = [\,A_\F;\,G_{\Sset\F}\,]$, so each
step is still one solve against $\Sig_{\F\F}$ (through the operator interface) feeding
an $(m+|\Sset|)$-dimensional Schur complement. The box bounds stay a separate per-variable
active set, so the box-only problem pays nothing for the generalisation.

The stacked solve \eqref{eq:schur} returns an affine multiplier
$\eta_j(\lambda) = \eta_j^{\alpha} + \lambda\,\eta_j^{\beta}$ for each active row, and
the two row events mirror the box events:
\begin{itemize}
\item \textbf{Activate.} An inactive row with slack
$s_j(\lambda) = (g_j\trans\alpha - h_j) + \lambda\,g_j\trans\beta$ binds at
$\lambda_j = (h_j - g_j\trans\alpha)/(g_j\trans\beta)$, admissible when
$g_j\trans\beta < 0$ (the slack rises as $\lambda$ falls).
\item \textbf{Release.} An active row releases at
$\lambda_j = -\eta_j^{\alpha}/\eta_j^{\beta}$, where its multiplier reaches zero,
admissible when $\eta_j^{\beta} > 0$.
\end{itemize}
These fill the trailing $p$ rows of $L$, so a group cap or factor-exposure band is
traced along the whole frontier (\S\ref{sec:exactness}). Each active row adds a column
to the solve against $\Sig_{\F\F}$ and a row to the Schur complement; with many active
rows that term, not the covariance solve, dominates a step. A row that is tight but a combination of the active rows on the
free set is supported: at the first vertex only an independent subset of the tight rows
is made active (\S\ref{sec:first}), and during the trace such a row has zero slope,
$g_j\trans\beta = 0$ (since $C_\F\beta_\F = 0$), so its activation event does not fire.
The active rows themselves must stay linearly independent on the free set, or the Schur
complement is singular, as when several caps bind at once in the same direction. Before
each solve the rows are tested, and a dependent set is refused with a
\code{DegenerateProblemError} at that $\lambda$.

\subsection{Gross-exposure caps}
\label{sec:leverage}

A cap $\lVert w\rVert_1 \le c$ on gross exposure (under a full budget, $c = 1.3$ is a
130/30 book) is polyhedral but not a single row of $G$. The code lifts it. Each asset
whose box straddles zero is split into a long and a short leg, $w_i = x_i^+ - x_i^-$
with $x_i^+\in[0,\ub_i]$ and $x_i^-\in[0,-\lb_i]$; a long-only or short-only asset keeps one
leg, so a long-only problem is not enlarged. With $w = Px$ for the signed
leg-to-asset map, the cap becomes the one row $\mathbf{1}\trans x \le c$ appended to
$GP$, the other constraints carry over as $AP$ and $GP$, and the mean becomes
$P\trans\mu$. The lifted program is an instance of \eqref{eq:qp} with the same
$\lambda$, traced by the loop above and mapped back through $P$ (validated in
\S\ref{sec:exactness}).

The lifted covariance $P\trans\Sig P$ is singular, since raising both legs of one
asset leaves $w$ unchanged. The loop never meets that singularity: both legs of an asset
are never free together, so every free block is a signed principal submatrix of $\Sig$,
served through the operator interface without forming $P\trans\Sig P$, and a factor
backend keeps its Woodbury solve. Appendix~\ref{sec:liftdetail} gives the argument and
two edge cases of the first vertex.

\subsection{Anti-cycling and termination}
\label{sec:bland}

On degenerate problems (tied returns, exchangeable assets, several constraints active at
once) several events can share $\lambda^{\star}$, and a naive rule can cycle. Among all
events within tolerance of $\lambda^{\star}$ the implementation picks the lowest
coordinate index, and within it the lowest event column, and flips that coordinate
alone. Call a partition (free set, blocked bounds, active rows $\Sset$) \emph{valid} at
$\lambda$ when its segment respects the bounds and the inactive rows there and its
multipliers have the right signs.

\begin{proposition}\label{prop:termination}
Let $\Sig \succ 0$, and suppose that at every turning point the constraints holding
with equality (the rows of $A$, the rows of $G$ with $g_j\trans w = h_j$, and the
bounds on which a weight sits) have linearly independent gradients. Then, in exact
arithmetic, the walk never returns to a partition it has left, and it reaches
$\lambda = 0$ after finitely many turning points.
\end{proposition}

\begin{proof}[Outline]
A partition stays valid while finitely many conditions, each affine in $\lambda$,
hold; once one fails below a turning point it fails for every smaller $\lambda$, so a
partition never recurs at a later $\lambda$, and without ties every step uses a new one.
At a tie, the slope of the next segment solves a strictly convex QP whose KKT system is a
linear complementarity problem with a positive-definite matrix, and the implementation's
lowest-index single flip is Murty's least-index principal pivot on it, which is finite.
Appendix~\ref{sec:proof} gives the full argument, including the correspondence between
the event types of the code and the complementary pairs of the problem.
\end{proof}

The proposition holds in exact arithmetic and under its hypotheses; what it adds is that
the walk never revisits a partition. Two weaker statements hold without them. Because
$\Sig \succ 0$ the optimiser $w(\lambda)$ is unique, and for a fixed set of active
constraints it is affine in $\lambda$ on the interval where that set is optimal; there
are finitely many such sets, so the frontier itself has finitely many turning points
whether or not the constraints are independent. And an iteration cap of $100(n+p+1)$
stops every trace: a trace that reaches it raises an error and never returns a partial
frontier, and no trace in Appendix~\ref{sec:degeneracy} used more than $1.5\%$ of it.
In exact arithmetic, then, the one thing left unproved outside the hypotheses is that
the walk leaves every tie without revisiting a partition. This concerns a degenerate maximum-return vertex,
resolved from a vertex of its optimal dual set, or, under tied maximum returns, from the
minimiser of the quadratic term on the tied face (\S\ref{sec:first}); a singular $\Sig$
that the active rows keep positive definite on their null space, outside
$\Sig\succ0$ (Appendix~\ref{sec:degeneracy}); the lifted problem
of \S\ref{sec:leverage}, whose singular covariance masks the competing leg events; and
ties with dependent constraints, where the matrix $M$ of Appendix~\ref{sec:proof} is
only positive semidefinite and the least-index result no longer applies (a partition
with dependent active rows is refused, \S\ref{sec:ineq}). Appendix~\ref{sec:degeneracy}
lists which test families satisfy the hypotheses. In floating point the tie window and
the slope floor make the computed signs agree with the exact ones, which is not proved
either. At run time every turning point is validated against the bounds, $A w = b$ and
$G w \le h$ before it is accepted, and the KKT certificate of \S\ref{sec:exactness}
checks the traces of \S\ref{sec:experiment} and Appendix~\ref{sec:degeneracy} for
optimality.

When no admissible event has a positive ratio the walk has reached the minimum-variance
portfolio and appends a final turning point $w = \alpha$ at $\lambda = 0$.

\begin{center}
\fbox{\begin{minipage}{0.93\textwidth}
\textbf{Algorithm 1.}~Critical Line Algorithm (as implemented in \pkg{cvxcla}).

\smallskip
\textit{Input:} $\mu$, $\Sig$, $\lb,\ub$, $A,b$, $G,h$, tolerance $\tau$.
\begin{enumerate}\setlength{\itemsep}{1pt}
\item Compute the maximum-return vertex (\S\ref{sec:first}); record its free set and
      active rows $\Sset$ with $\lambda \gets \infty$.
\item \textbf{While} $\lambda > 0$:
  \begin{enumerate}\setlength{\itemsep}{1pt}
  \item Fix $w_\B$ at the active bounds and stack $C_\F = [\,A_\F;\,G_{\Sset\F}\,]$.
  \item Solve against $\Sig_{\F\F}$ for $Y$, $z^{\alpha}$, $z^{\beta}$; solve the Schur
        system \eqref{eq:schur}; back-substitute \eqref{eq:freesoln} for
        $\alpha,\beta$; form $\gamma,\delta$ \eqref{eq:gamma}.
  \item Assemble $L$ (\S\ref{sec:events}, \S\ref{sec:ineq}) under the slope floor and
        $\lambda$ window; $\lambda^{\star} \gets \max L$; \textbf{if}
        $\lambda^{\star} \le 0$ \textbf{break}.
  \item Pick the event by the least-index rule (\S\ref{sec:bland}), flip that
        coordinate's status, set
        $\lambda \gets \lambda^{\star}$, append $w = \alpha + \lambda^{\star}\beta$.
  \end{enumerate}
\item Append $w = \alpha$ at $\lambda = 0$ and \textbf{return} the turning points.
\end{enumerate}
\end{minipage}}
\end{center}
